\chapter{Matrices et graphes}\label{ch:g12:matrix}

Une \omterm{def:g12:matrix:matrix}{matrice} est un tableau rectangulaire de nombres, additionné et multiplié
selon des règles conçues pour que l'algèbre matricielle représente la
composition des applications linéaires. Les \omterm{def:g12:matrix:matrix}{matrices} résolvent les \omterm{def:g10:lines:system}{systèmes
linéaires}, pilotent les \omterm{def:g12:seq:sequence}{suites} récurrentes couplées, et comptent les
chemins dans les réseaux --- les mathématiques derrière les moteurs de
recherche et les algorithmes de plus court chemin.

\section{Algèbre matricielle}

\begin{definition}[Matrice]\label{def:g12:matrix:matrix}
Une \emph{matrice $m \times n$}\index{matrice} est un tableau de \omterm{def:g10:numbers:sets}{nombres
réels} à $m$ lignes et $n$ colonnes~:
$A = (a_{ij})$, où $a_{ij}$ est l'entrée de la ligne $i$, colonne $j$.
Deux matrices de même taille s'additionnent entrée par entrée, et
$\lambda A = (\lambda a_{ij})$.
\end{definition}

\begin{definition}[Produit matriciel]\label{def:g12:matrix:product}
Soit $A$ de taille $m \times n$ et $B$ de taille $n \times p$. Le produit
$AB$ est la \omterm{def:g12:matrix:matrix}{matrice} $m \times p$ dont l'entrée $(i,j)$ est
\[
(AB)_{ij} = \sum_{k=1}^{n} a_{ik} b_{kj}
\]
(la règle «~ligne $i$ de $A$ multipliée par colonne $j$ de $B$~»).
\end{definition}

\begin{example}
$\begin{pmatrix} 1 & 2\\ 3 & 4\end{pmatrix}
\begin{pmatrix} 0 & 1\\ 1 & 1\end{pmatrix}
= \begin{pmatrix} 2 & 3\\ 4 & 7\end{pmatrix}$,
\quad tandis que \quad
$\begin{pmatrix} 0 & 1\\ 1 & 1\end{pmatrix}
\begin{pmatrix} 1 & 2\\ 3 & 4\end{pmatrix}
= \begin{pmatrix} 3 & 4\\ 4 & 6\end{pmatrix}$~:
\emph{la multiplication matricielle n'est pas commutative}.
\end{example}

\begin{proposition}[Règles de l'algèbre matricielle]\label{prop:g12:matrix:rules}
Dès que les tailles rendent les produits significatifs~:
\[
(AB)C = A(BC), \qquad A(B + C) = AB + AC, \qquad (A+B)C = AC + BC,
\]
et la \emph{\omterm{def:g12:matrix:matrix}{matrice} identité} $I_n$ (uns sur la diagonale, zéros ailleurs)
satisfait $I_m A = A I_n = A$ pour $A$ de taille $m \times n$.
\end{proposition}

\begin{proof}
Toutes se vérifient entrée par entrée à partir de
\cref{def:g12:matrix:product}~; l'associativité, la seule non triviale,
revient à échanger deux sommes finies~:
$\bigl((AB)C\bigr)_{ij} = \sum_l \left(\sum_k a_{ik}b_{kl}\right) c_{lj}
= \sum_k a_{ik} \left(\sum_l b_{kl} c_{lj}\right)
= \bigl(A(BC)\bigr)_{ij}$.
\end{proof}

\begin{definition}[Inverse]\label{def:g12:matrix:inverse}
Une \omterm{def:g12:matrix:matrix}{matrice} carrée $A$ de taille $n$ est \emph{inversible}\index{matrice!inverse}
s'il existe une \omterm{def:g12:matrix:matrix}{matrice} $B$ avec $AB = BA = I_n$~; $B$ est alors unique,
notée $A^{-1}$.
\end{definition}

\begin{proposition}[Inverse d'une matrice $2\times2$]\label{prop:g12:matrix:inverse2x2}
Soit $A = \begin{pmatrix} a & b\\ c & d\end{pmatrix}$ et
$\det A = ad - bc$ (le \emph{déterminant}\index{déterminant}). Alors $A$ est
\omterm{def:g12:matrix:inverse}{inversible} si et seulement si $\det A \neq 0$, auquel cas
\[
A^{-1} = \frac{1}{ad - bc}\begin{pmatrix} d & -b\\ -c & a\end{pmatrix}.
\]
\end{proposition}

\begin{proof}
Un calcul donne
$A \begin{pmatrix} d & -b\\ -c & a\end{pmatrix}
= \begin{pmatrix} d & -b\\ -c & a\end{pmatrix} A = (ad - bc) I_2$~; si
$ad - bc \neq 0$, diviser. Réciproquement, si $ad - bc = 0$, les colonnes de
$A$ sont proportionnelles, et de même les colonnes de $AB$ pour toute $B$~;
mais les colonnes de $I_2$ ne le sont pas, donc aucune $B$ ne peut vérifier
$AB = I_2$.
\end{proof}

\begin{method}[Systèmes linéaires]\label{met:g12:matrix:systems}
Le système
$\begin{cases} ax + by = e\\ cx + dy = f \end{cases}$
est l'\omterm{def:g10:algebra:equation}{équation} matricielle $AX = Y$ avec
$X = \begin{pmatrix} x \\ y\end{pmatrix}$,
$Y = \begin{pmatrix} e \\ f\end{pmatrix}$. Si $\det A \neq 0$, sa solution
unique est $X = A^{-1}Y$. Le même formalisme traite $n$ \omterm{def:g10:algebra:equation}{équations} à $n$
inconnues.
\end{method}

\section{Puissances de matrices et suites récurrentes}

\begin{definition}\label{def:g12:matrix:power}
Pour une \omterm{def:g12:matrix:matrix}{matrice} carrée $A$ et $k \in \N$, $A^k = A \times \dots \times A$
($k$ facteurs), avec $A^0 = I$.
\end{definition}

\begin{method}[Cas diagonal-plus-nilpotent et diagonalisable]\label{met:g12:matrix:powers}
Deux façons standards de calculer $A^k$~:
\begin{itemize}
  \item Si $A = \lambda I + N$ avec $N^2 = 0$, la formule du binôme
        (valable ici car $I$ et $N$ commutent) se réduit à deux termes~:
        $A^k = \lambda^k I + k \lambda^{k-1} N$.
  \item Si l'on trouve une $P$ \omterm{def:g12:matrix:inverse}{inversible} avec $A = PDP^{-1}$ et $D$
        diagonale, alors $A^k = P D^k P^{-1}$, et $D^k$ se calcule entrée
        par entrée. (Trouver une telle $P$ de façon systématique est la
        théorie de la \emph{diagonalisation}, développée à l'université~; à
        ce niveau $P$ est donnée.)
\end{itemize}
\end{method}

\begin{example}[Suites couplées]\label{ex:g12:matrix:coupled}
Soient $u_{n+1} = 3u_n + v_n$ et $v_{n+1} = u_n + 3v_n$. En posant
$X_n = \begin{pmatrix} u_n\\ v_n \end{pmatrix}$ et
$A = \begin{pmatrix} 3 & 1\\ 1 & 3\end{pmatrix}$, on obtient
$X_{n+1} = AX_n$, donc $X_n = A^n X_0$. Les \omterm{def:g12:seq:sequence}{suites} auxiliaires
$s_n = u_n + v_n$ et $d_n = u_n - v_n$ vérifient $s_{n+1} = 4s_n$ et
$d_{n+1} = 2d_n$, donc $s_n = 4^n s_0$, $d_n = 2^n d_0$ et
\[
u_n = \frac{4^n(u_0+v_0) + 2^n(u_0-v_0)}{2}, \qquad
v_n = \frac{4^n(u_0+v_0) - 2^n(u_0-v_0)}{2}.
\]
(En coulisse~: $(1,1)$ et $(1,-1)$ sont des directions propres de $A$.)
\end{example}

\section{Graphes et chemins}

\begin{definition}[Graphe, matrice d'adjacence]\label{def:g12:matrix:graph}
Un \emph{graphe}\index{graphe} est constitué de sommets $1, 2, \dots, n$ et
d'arêtes joignant certaines paires de sommets (couples ordonnés pour un
graphe \emph{orienté}). Sa \emph{matrice d'adjacence}\index{matrice d'adjacence}
est la \omterm{def:g12:matrix:matrix}{matrice} $n \times n$ $M$ avec $m_{ij} = 1$ s'il y a une arête de $i$
vers $j$, et $0$ sinon. Un \emph{chemin} de longueur $k$ de $i$ vers $j$ est
une \omterm{def:g12:seq:sequence}{suite} de $k$ arêtes consécutives menant de $i$ à $j$.
\end{definition}

\begin{omfigure}
\begin{tikzpicture}[baseline=(current bounding box.center)]
  \node[circle, draw, inner sep=2.5pt] (v1) at (0,0) {$1$};
  \node[circle, draw, inner sep=2.5pt] (v2) at (2.6,1.2) {$2$};
  \node[circle, draw, inner sep=2.5pt] (v3) at (2.6,-1.2) {$3$};
  \draw[-Stealth, omDef, thick] (v1) -- (v2);
  \draw[-Stealth, omDef, thick] (v2) -- (v3);
  \draw[-Stealth, omDef, thick] (v1) to[bend right=18] (v3);
  \draw[-Stealth, omDef, thick] (v3) to[bend right=18] (v1);
\end{tikzpicture}
\qquad\qquad
$M = \begin{pmatrix} 0 & 1 & 1\\ 0 & 0 & 1\\ 1 & 0 & 0 \end{pmatrix}$

{\small Un \omterm{def:g12:matrix:graph}{graphe} orienté et sa \omterm{def:g12:matrix:graph}{matrice d'adjacence}
(\cref{exo:g12:matrix:6})~: $m_{ij} = 1$ exactement lorsqu'il y a une arête
de $i$ vers $j$.}
\end{omfigure}

\begin{theorem}[Dénombrement des chemins]\label{thm:g12:matrix:walks}
Le nombre de chemins de longueur $k$ du sommet $i$ au sommet $j$ est
l'entrée $(i,j)$ de $M^k$.
\end{theorem}

\begin{proof}
Récurrence sur $k$. Pour $k = 1$ c'est la définition de $M$. Supposer le
résultat pour $k$. Un chemin de longueur $k+1$ de $i$ à $j$ est un chemin
de longueur $k$ de $i$ vers un certain sommet $l$, suivi d'une arête de $l$
vers $j$~; par les principes d'addition et de multiplication, leur nombre
est
\[
\sum_{l=1}^{n} \bigl(M^k\bigr)_{il}\, m_{lj} = \bigl(M^{k+1}\bigr)_{ij}.
\qedhere
\]
\end{proof}

\begin{example}
Pour le \omterm{def:g12:matrix:graph}{graphe} triangle ($3$ sommets, toutes les paires jointes),
$M = \begin{pmatrix} 0&1&1\\ 1&0&1\\ 1&1&0\end{pmatrix}$ et
$M^2 = \begin{pmatrix} 2&1&1\\ 1&2&1\\ 1&1&2\end{pmatrix}$~: depuis chaque
sommet il y a $2$ chemins de longueur $2$ vers soi-même (via l'un ou l'autre
voisin) et $1$ vers chaque autre sommet.
\end{example}

\section{Exercices}

\begin{exercise}[$\star$]\label{exo:g12:matrix:1}
Soient $A = \begin{pmatrix} 1 & 2\\ 0 & 1 \end{pmatrix}$ et
$B = \begin{pmatrix} 2 & 0\\ 1 & 1 \end{pmatrix}$. Calculer $A + B$, $AB$,
$BA$ et $A^2$.
\end{exercise}

\begin{exercise}[$\star$]\label{exo:g12:matrix:2}
Déterminer si les \omterm{def:g12:matrix:matrix}{matrices} suivantes sont \omterm{def:g12:matrix:inverse}{inversibles}, et calculer les
inverses lorsqu'ils existent~:
\[
A = \begin{pmatrix} 2 & 5\\ 1 & 3\end{pmatrix}, \qquad
B = \begin{pmatrix} 3 & 6\\ 2 & 4\end{pmatrix}.
\]
\end{exercise}

\begin{exercise}[$\star$]\label{exo:g12:matrix:3}
Résoudre par inversion matricielle le système
$\begin{cases} 2x + 5y = 1\\ x + 3y = 2 . \end{cases}$
\end{exercise}

\begin{exercise}[$\star\star$]\label{exo:g12:matrix:4}
Soit $A = \begin{pmatrix} 2 & 1\\ 0 & 2\end{pmatrix} = 2I + N$ avec
$N = \begin{pmatrix} 0 & 1\\ 0 & 0 \end{pmatrix}$.
\begin{enumerate}
  \item Vérifier que $N^2 = 0$ et que $I$ et $N$ commutent.
  \item En déduire $A^k$ pour tout $k \in \N$ et vérifier la formule pour
        $k=2$ par calcul direct.
\end{enumerate}
\end{exercise}

\begin{exercise}[$\star\star$]\label{exo:g12:matrix:5}
Soient $A = \begin{pmatrix} 0 & 1\\ 1 & 1\end{pmatrix}$ et $F_n$ la \omterm{def:g12:seq:sequence}{suite}
de Fibonacci ($F_0 = 0$, $F_1 = 1$, $F_{n+2} = F_{n+1} + F_n$). Montrer par
récurrence que pour $n \geq 1$,
\[
A^n = \begin{pmatrix} F_{n-1} & F_n\\ F_n & F_{n+1}\end{pmatrix},
\]
et en déduire l'identité $F_{n+1}F_{n-1} - F_n^2 = (-1)^n$.
(Indication~: les \omterm{def:g11:vect:det}{déterminants} se multiplient~: $\det(MN) = \det M \det N$,
ce que l'on peut vérifier pour les \omterm{def:g12:matrix:matrix}{matrices} $2\times2$.)
\end{exercise}

\begin{exercise}[$\star\star$]\label{exo:g12:matrix:6}
Un \omterm{def:g12:matrix:graph}{graphe} orienté sur les sommets $\{1, 2, 3\}$ a les arêtes
$1\to2$, $2\to3$, $3\to1$ et $1\to3$.
\begin{enumerate}
  \item Écrire la \omterm{def:g12:matrix:graph}{matrice d'adjacence} $M$ et calculer $M^2$ et $M^3$.
  \item Combien de chemins de longueur $3$ vont de $1$ à $1$~? Les lister.
\end{enumerate}
\end{exercise}

\begin{exercise}[$\star\star$]\label{exo:g12:matrix:7}
Une société d'autopartage déplace des véhicules entre deux villes $A$ et
$B$. Chaque semaine, $80\%$ des voitures en $A$ restent en $A$ et $20\%$
passent en $B$~; $30\%$ des voitures en $B$ passent en $A$ et $70\%$ restent.
Soient $a_n, b_n$ les proportions de la flotte dans chaque ville.
\begin{omfigure}
\begin{tikzpicture}
  \node[circle, draw, inner sep=3pt] (A) at (0,0) {$A$};
  \node[circle, draw, inner sep=3pt] (B) at (3.2,0) {$B$};
  \draw[-Stealth, omDef, thick] (A) to[bend left=25]
    node[above, font=\small] {$0.2$} (B);
  \draw[-Stealth, omDef, thick] (B) to[bend left=25]
    node[below, font=\small] {$0.3$} (A);
  \draw[-Stealth, omDef, thick] (A) to[loop left]
    node[left, font=\small] {$0.8$} (A);
  \draw[-Stealth, omDef, thick] (B) to[loop right]
    node[right, font=\small] {$0.7$} (B);
\end{tikzpicture}
\end{omfigure}
\begin{enumerate}
  \item Écrire $X_{n+1} = MX_n$ avec
        $X_n = \begin{pmatrix} a_n\\ b_n\end{pmatrix}$ et identifier $M$.
  \item Trouver les proportions d'équilibre (résoudre $MX = X$ avec
        $a + b = 1$).
  \item Montrer que $c_n = a_n - 0.6$ vérifie $c_{n+1} = 0.5\,c_n$, et
        conclure que la répartition de la flotte \omterm{def:g12:seq:limit}{converge} vers l'équilibre.
\end{enumerate}
\end{exercise}

\begin{exercise}[$\star\star\star$]\label{exo:g12:matrix:8}
Soient $A = \begin{pmatrix} 3 & 1\\ 1 & 3 \end{pmatrix}$,
$P = \begin{pmatrix} 1 & 1\\ 1 & -1 \end{pmatrix}$.
\begin{enumerate}
  \item Calculer $P^{-1}$, puis $D = P^{-1}AP$, et vérifier que $D$ est
        diagonale.
  \item En déduire une formule fermée pour $A^n$ et comparer avec
        \cref{ex:g12:matrix:coupled}.
\end{enumerate}
\end{exercise}

\section{Problème~: la matrice qui connaît Fibonacci (et la météo)}

\begin{problem}[{Devoir du week-end --- une seule matrice
$2 \times 2$ porte tout Fibonacci, une matrice de Markov prévoit la
météo à long terme, et un vecteur propre vaut un milliard}]
\label{pb:g12:matrix:1}
Une \omterm{def:g12:matrix:matrix}{matrice} est une machine qui avale un état et renvoie le
suivant --- et ses \emph{puissances} contiennent donc des avenirs
entiers. Ce problème s'ouvre sur l'étonnante \omterm{def:g12:matrix:matrix}{matrice} dont les
puissances égrènent les nombres de Fibonacci (et démontrent leurs
identités en une ligne chacune), fait ensuite tourner la météo comme
une chaîne de Markov jusqu'à son état stationnaire, et se referme
sur le \omterm{def:g11:vect:vector}{vecteur} propre sur lequel un moteur de recherche a été bâti
(\cref{thm:g12:matrix:walks}, \cref{met:g12:matrix:powers}).

\medskip
\noindent\textbf{Partie I --- Aisance.}
\begin{enumerate}
  \item Avec $A = \begin{pmatrix} 1 & 2\\ 3 & 4\end{pmatrix}$ et
        $B = \begin{pmatrix} 0 & 1\\ 1 & 0\end{pmatrix}$~:
        calculer $AB$ et $BA$. Verdict sur la commutativité~?
  \item Inverser $\begin{pmatrix} 2 & 1\\ 5 & 3\end{pmatrix}$
        (\cref{prop:g12:matrix:inverse2x2}) et se servir de
        l'inverse pour résoudre $2x + y = 4$, $5x + 3y = 7$.
  \item Soit $N = \begin{pmatrix} 0 & 1\\ 0 & 0\end{pmatrix}$~:
        calculer $N^2$, et en déduire $(I + N)^n = I + nN$ pour
        tout $n$.
  \item Le \omterm{def:g12:matrix:graph}{graphe} triangle (trois sommets, toutes les paires
        reliées)~: écrire sa \omterm{def:g12:matrix:graph}{matrice d'adjacence} $A$, calculer
        $A^3$, et interpréter les coefficients diagonaux
        (\cref{thm:g12:matrix:walks}).
  \item Pour $D = \begin{pmatrix} 2 & 0\\ 0 & \frac12
        \end{pmatrix}$~: donner $D^n$ et son comportement quand
        $n \to \infty$.
\end{enumerate}

\medskip
\noindent\textbf{Partie II --- La \omterm{def:g12:matrix:matrix}{matrice} de Fibonacci.}
Posons $F = \begin{pmatrix} 1 & 1\\ 1 & 0\end{pmatrix}$ et notons
$F_1 = F_2 = 1, F_3 = 2, \dots$ les nombres de Fibonacci du
\cref{pb:g11:seq:1}.
\begin{enumerate}[resume]
  \item Calculer $F^2$, $F^3$, $F^4$ et conjecturer la forme
        générale de $F^n$ en \omterm{def:g11:func:function}{fonction} des nombres de Fibonacci.
  \item Démontrer la conjecture
        $F^n = \begin{pmatrix} F_{n+1} & F_n\\ F_n & F_{n-1}
        \end{pmatrix}$ par récurrence.
  \item Prendre le \omterm{def:g11:vect:det}{déterminant} des deux membres (le \omterm{def:g11:vect:det}{déterminant}
        d'un produit est le produit des \omterm{def:g11:vect:det}{déterminants} --- le
        vérifier sur des \omterm{def:g12:matrix:matrix}{matrices} $2 \times 2$ si vous ne l'avez
        jamais vu)~: en déduire l'\emph{identité de Cassini}
        $F_{n+1}F_{n-1} - F_n^2 = (-1)^n$ --- le moteur du carré
        évanoui, démontré en une ligne.
  \item À partir de $F^{m+n} = F^m F^n$, lire les coefficients en
        haut à droite et en déduire la \emph{formule d'addition}
        \[
        F_{m+n} = F_{m+1} F_n + F_m F_{n-1} .
        \]
        La vérifier pour $m = n = 3$.
  \item Déduire de la formule d'addition (récurrence sur $k$) que
        $F_n$ divise $F_{kn}$, et le vérifier sur $F_3 \mid F_6$ et
        $F_3 \mid F_9$.
  \item Pour calculer $F_{100}$, nul besoin de multiplier $100$
        \omterm{def:g12:matrix:matrix}{matrices}~: élever au carré de façon répétée
        ($F^2, F^4, F^8, \dots$) puis combiner. Combien de
        multiplications de \omterm{def:g12:matrix:matrix}{matrices} suffisent, et de quelle antique
        astuce de multiplication du volume précédent s'agit-il,
        promue aux \omterm{def:g12:matrix:matrix}{matrices}~?
\end{enumerate}

\medskip
\noindent\textbf{Partie III --- La machine à météo.}
Dans une certaine ville~: après une journée ensoleillée, la suivante
est ensoleillée avec la \omterm{def:g10:proba:distribution}{probabilité} $0.8$~; après une journée
pluvieuse, elle est ensoleillée avec la \omterm{def:g10:proba:distribution}{probabilité} $0.4$. On code
la \omterm{def:g11:prob:rv}{loi} du jour par une colonne
$\binom{p_{\text{soleil}}}{p_{\text{pluie}}}$ et l'évolution par
\[
M = \begin{pmatrix} 0.8 & 0.4\\ 0.2 & 0.6 \end{pmatrix}.
\]
\begin{enumerate}[resume]
  \item Vérifier que chaque colonne de $M$ a pour somme $1$, et
        dire pourquoi toute machine à météo doit posséder cette
        propriété.
  \item Aujourd'hui il fait soleil. Calculer la prévision pour
        demain, puis pour après-demain.
  \item Déterminer l'\emph{état stationnaire}~: la \omterm{def:g11:prob:rv}{loi} $v$ telle
        que $Mv = v$ (et dont les coefficients ont pour somme $1$).
        Quelle fraction des jours est ensoleillée à long terme~?
  \item Partir d'une journée pluvieuse, $\binom01$, et appliquer
        $M$ quatre fois en suivant à chaque pas la distance à l'état
        stationnaire. Par quel facteur l'écart se réduit-il à chaque
        pas --- et de quel type de convergence s'agit-il~?
  \item Le PageRank en miniature~: trois pages, avec les liens
        $A \to B$, $A \to C$, $B \to C$, $C \to A$. Un internaute
        aléatoire suit un lien sortant uniformément au hasard.
        Écrire la \omterm{def:g12:matrix:matrix}{matrice} de transition, déterminer l'état
        stationnaire, et classer les pages.
  \item Interpréter le classement~: pourquoi $C$ obtient-il un
        score aussi élevé que $A$ alors qu'il reçoit des liens de
        moins de pages --- que mesure au juste l'état
        stationnaire~? (Le vrai PageRank ajoute un facteur
        d'amortissement pour les impasses et les sauts~; l'idée du
        \omterm{def:g11:vect:vector}{vecteur} propre est exactement celle-ci.)
\end{enumerate}

\medskip
\noindent\textbf{Partie IV --- Les dividendes de la diagonale.}
\begin{enumerate}[resume]
  \item Deux quantités couplées obéissent à
        $u_{n+1} = 3u_n + v_n$, $v_{n+1} = u_n + 3v_n$,
        c'est-à-dire à la \omterm{def:g12:matrix:matrix}{matrice} $A$ de l'\cref{exo:g12:matrix:8}.
        À l'aide de la diagonalisation de cet exercice
        ($D = \operatorname{diag}(4, 2)$), donner la formule close
        de $u_n$ pour $u_0 = 1$, $v_0 = 0$, et la confronter au
        calcul direct pour $n = 1, 2, 3$.
  \item En une ou deux phrases~: que \emph{fait} la
        diagonalisation à un système couplé --- et en quel sens
        l'état stationnaire de Markov de la question 14 est-il lui
        aussi une histoire de \omterm{def:g11:vect:vector}{vecteur} propre~?
  \item Pour finir --- les trois visages de la \omterm{def:g12:matrix:matrix}{matrice} ce week-end~:
        la tenue de comptes (systèmes et inverses), le dénombrement
        (chemins et liens comptés par les puissances) et l'évolution
        (Fibonacci, la météo, le web --- des avenirs lus sur les
        directions propres). Une phrase pour chacun, plus l'annonce~:
        l'algèbre linéaire des volumes universitaires fait de chacun
        de ces visages une théorie.
\end{enumerate}
\end{problem}
